\documentclass{hw_template}
\usepackage{xcolor}
\usepackage{graphicx}

% STUDENTS: Please fill these out with your information
\newcommand{\name}{Dylan Andres}
\newcommand{\netid}{dra77}
\newcommand{\collaborators}{}
% / STUDENT

\newcommand{\hwnum}{1}
\newcommand{\duedate}{September 23, 11:59pm ET}
\renewcommand{\title}{Stable Matching, Greedy, Divide and Conquer}
\newcommand{\io}{\textbf{Code input and output format.} }
\newcommand{\submission}{\textbf{Submission.}}
\newcommand{\floor}[1]{\lfloor #1 \rfloor}

% \newtheorem{claim}{Claim}

\begin{document}

%%%%%%%%%%%%%%%%%%%%%%%%%%%%%%%%%%%%%%%%%%%%%%%%%%%%%%%%%%%%%%%%%%%%%%%%%%%%%%%%
% Problem 1: G-S algorithm

\begin{problem}
In this problem, you will compute the maximum sum of non-adjacent elements,
$\operatorname{maxsum}(x)
= \max_{S \in \mathcal{S}} \; \sum_{s \in S} x_s$
where no two elements in \(S\) are adjacent, in two settings: a list and a tree. The definition of
“adjacent” will be specific to each setting. Additionally, for each setting you must first state
the subproblem precisely and then implement the full dynamic programming algorithm.\\


(a) (\textit{2 points})
Given a sequence of integers $x = (x_1, x_2, \ldots, x_n)$,
your goal is to compute the largest sum of non-adjacent elements of \(x\), denoted as
\(\operatorname{maxsum\_list}(x)\).
Adjacent elements of the sequence are next to each other: \(x_i\) and \(x_{i+1}\). \\\\
Identify the dynamic programming subproblem and write its solution here.\\

\textbf{The DP subproblem is: In order to maximize the sum of non-adjacent integers in a list, should we 
take the current integer plus the total sum of the integer two places before it? Or should we
skip the current integer and take the total sum of the integer one place before it?}\\

\textbf{The solution involves creating a DP array of size n+1 (where n is the length of input list) so 
that it can be 1-indexed. The base cases are dp[0] = 0 and dp[1] = x[0] (where dp is the DP array and x is the input list).
This is because an empty list has a max sum of 0 and a list with one element has a max sum of that element.
Then, for each index i from 2 to n, we compute dp[i] as the maximum of  dp[i-2] + x[i-1] 
(taking the current integer plus the total sum of the integer two places before it), and dp[i-1] 
(skipping current element and taking the total sum of the integer one place before it).
Finally, we return dp[n] as the max sum.}\\\\

(b) (\textit{2 points})
Given a tree with nodes that contain integers, your goal is to compute the largest sum
of non-adjacent elements. Trees are represented as a set of vertices
$v_i \in V$ and edges $(v_i, v_j) \in E \subseteq V \times V$.
There are no loops and every vertex has at most a single parent.
Adjacent elements of the tree are those that have an edge between nodes, i.e. $(v_i, v_j) \in E$.
For convenience, define the children of a node to be
$c(v_i) = { v_j \mid (v_i, v_j) \in E }$\\\\
Identify the dynamic programming subproblem and write its solution here.\\\\

\textbf{The DP subproblem is: In order to maximum the sum of non-adjacent nodes in a tree, should we
take the current node's value plus the total sum of its sub-trees (excluding its direct child nodes) or should we 
skip the current node's value and take the maximum sum from all of its sub-trees}\\

\textbf{The solution starts by evaluating two edge cases, a tree of size 0 has a maximum of 0, and a tree of size 1 has
a maximum of the value of the single node. Next, we find the root of tree to start the heuristic. Since a node can 
have mulitple children, the tree has to be traversed recursively starting with the root. At each step of the recursive sequence,
we compute for each child:}\\
\textbf{1. The sum of taking the current node's value plus the total sum of its sub-trees (excluding its direct child nodes)}\\
\textbf{2. The sum of skipping the current node's value and taking the maximum sum from all of its sub-trees}\\
\textbf{These sums are returned up the recursive call stack back to the root. Finally, we return the maximum of the two sums.}\\\\

\end{problem}

\newpage
\begin{problem}
Given $n$ points in a 2D plane, your goal is to find the two points that are closest
together according to the Euclidean distance
$d(x, y) = \sqrt{(x_i - y_i)^2 + (x_j - y_j)^2}$.


(d) (\textit{5 points})
Perform an empirical analysis of the three approaches, plotting their runtime
performance for a series of problem sizes. \\\\

\textbf{I ran a benchmark for each implementation on several input sizes as large as 5000, then 
plotted the median runtime. The resulting graph is very consistent with the asymptotic bounds of each algorithm. 
The brute-force algorithm (slowest) is clearly $O(n^2)$ as evident by the quadratic curve. The divide \& conquer looks 
linear at first, but there is a clear "elbow" around n = 2000, which is consistent with the $O(n \log n)$ bound. Lastly,
the randomized algorithm is the fastest and appears linear, which is consistent with its expected $O(n)$ time complexity. 
This is not as obvious, so I also plotted the randomized algorithm own its own:}\\


OS: MacOS Sequoia 15.6.1\\
CPU: Apple M1 Pro 8-core\\
Memory: 16 GB\\
**ChatGPT was used to generate example test cases for this problem\\
\\

\end{problem}

\newpage
\begin{problem}
\textbf{Skip lists} are randomized data structures that allow efficient (in expectation) 
insertion, search, and deletion operations in a sorted list. We assume that list elements 
are always integers.

You will implement skip lists in this problem, as well as a linked list baseline. 
We supply separate \texttt{Node} classes in each file \texttt{problem\_3/*}.py, 
which will be the backbone of each implementation.

Both implementations will have the following interface:
\begin{itemize}
    \item \textbf{insert(list, x: int)} $\rightarrow$ \texttt{None}: Add a new integer to our list.
    \item \textbf{search(list, x: int)} $\rightarrow$ \texttt{Node}: Return the node in our list that has the value $x$.
    \item \textbf{delete(list, x: int)} $\rightarrow$ \texttt{None}: Delete the node in our list that has value $x$.
\end{itemize}

For simplicity, we assume that no repeated elements will be inserted into the list.

Please read the following resource on skip lists for more helpful information.


(a) (\textit{5 points})
As a baseline, implement a linked list class. Give the runtime complexity of 
search, insertion, and deletion. The linked list must always be in sorted order, and all 
elements will be integers. Include your implementation in 
\texttt{problem\_3/p3\_a.py}.\\\\
\textbf{Search: O(n). The worst case is when the element to search does not exist, or when the element is the largest of all existing elements. All n nodes must be traversed.}\\\\
\textbf{Insertion: O(n). The worse case is when the element to insert is larger than all existing elements. All n nodes must be traversed.}\\\\
\textbf{Deletion: O(n). The worst case is when the element to delete does not exist or when the element is the largest of all existing elements. All n nodes must be traversed.}\\\\
\newpage

(b) (\textit{5 points})
Implement insertion for skip lists, which should take O(log n) expected time. All
list elements will be ints. Perform an empirical analysis of insertion in skip lists versus the
linked list baseline. Include your implementation in 
\texttt{problem\_3/p3\_a.py}.\\\\
\textbf{I ran a benchmark for both insert implementations on several input sizes as large as 5000. The resulting graph lines up with the 
expected asymptotic bound of each. LinkedList insertion is O(n), while SkipList insertion is O(logN). To make this more obvious, I plotted SkipList
on its own. The resulting graph has some noise since the algorithm is randomized. However, the general shape aligns with a logarithmic plot.}\\
\textbf{When timing the benchmark across N inputs, the graph also aligns with the expected asymptotic bounds.
LinkedList insertion is $O(n^2)$ since it is $O(n)$ across N operations.
SkipList insertion is $O(nlogN)$ since it is $O(logN)$ across N operations.}\\

\newpage

(c) (\textit{5 points})
Implement search for skip lists, which should take $O(\log n)$ expected time. Give
the expected and worst-case runtime complexity. All list elements will be integers. Perform
an \emph{empirical analysis} of search in skip lists versus the linked list baseline. Include your
implementation in \texttt{problem\_3/p3\_b.py}.\\\\
\textbf{I ran a benchmark for both search implementations on several input sizes as large as 5000. The resulting graph lines up with the 
expected asymptotic bound of each. LinkedList search is O(n), while SkipList search is O(logN). To make this more obvious, I plotted SkipList
on its own. The resulting graph has some noise since the algorithm is randomized. However, the general shape aligns with a logarithmic plot.}\\
\textbf{When timing the benchmark across N inputs, the graph also aligns with the expected asymptotic bounds.
LinkedList search is $O(n^2)$ since it is $O(n)$ across N operations.
SkipList search is $O(nlogN)$ since it is $O(logN)$ across N operations.}\\
\newpage

(d) (\textit{5 points})
Implement deletion for skip lists, which should take $O(\log n)$ expected time. All
list elements will be integers. Perform an \emph{empirical analysis} of deletion in skip lists versus the
linked list baseline. Include your implementation in \texttt{problem\_3/p3\_b.py}.\\\\
\textbf{I ran a benchmark for both delete implementations on several input sizes as large as 5000. The resulting graph lines up with the 
expected asymptotic bound of each. LinkedList deletion is O(n), while SkipList deletion is O(logN). To make this more obvious, I plotted SkipList
on its own. The resulting graph has some noise since the algorithm is randomized. However, the general shape aligns with a logarithmic plot.}\\
\textbf{When timing the benchmark across N inputs, the graph also aligns with the expected asymptotic bounds.
LinkedList deletion is $O(n^2)$ since it is $O(n)$ across N operations.
SkipList deletion is $O(nlogN)$ since it is $O(logN)$ across N operations.}\\

OS: MacOS Sequoia 15.6.1\\
CPU: Apple M1 Pro 8-core\\
Memory: 16 GB\\
**ChatGPT was used to generate example test cases for this problem\\
\\

\end{problem}

\end{document}
